\newpage
\section{Podsumowanie}

\subsection{Wnioski z przeprowadzonych badań}

Część badawcza (rozdział~\ref{sec:modele}) przebiegała jako sekwencja czterech
kolejnych odkryć, z~których każde ukierunkowywało następny eksperyment. Uzyskane
wyniki pozwoliły sformułować pięć wniosków, które łącznie potwierdzają, że
uruchomienie użytecznego asystenta głosowego opartego na małym modelu językowym
(SLM) na urządzeniu brzegowym klasy Raspberry~Pi~5 jest wykonalne.

Po pierwsze, stabilność termiczna platformy brzegowej okazała się warunkiem
koniecznym wiarygodnego pomiaru. Instalacja aktywnego chłodzenia sterowanego
GPIO zredukowała udział dławienia termicznego ze~100\% (Run~1) do~0\% (Run~2),
dzięki czemu ranking modeli przestał odzwierciedlać odporność na throttling,
a~zaczął mierzyć jakość semantyczną. Model docelowy \texttt{gemma3:4b} awansował
z~czwartego na pierwsze miejsce: \emph{composite score} wzrósł z~0{,}497 do~0{,}806
($+62{,}3\%$), a~mediana E2E skróciła się z~57\,584~ms do~31\,742~ms
($1{,}81\times$ szybciej) --- por.~tabele~\ref{tab:ranking_run1}
i~\ref{tab:ranking_run2}.

Po drugie, analiza statystyczna Run~3 wykazała, że o~wyniku decyduje przede
wszystkim flaga halucynacji: sama \texttt{hallucination\_flag} wyjaśnia
$R^2 = 0{,}9922$ wariancji \emph{composite score}, podczas gdy hiperparametry
inferencji zaledwie $R^2 = 0{,}2997$ (rysunek~\ref{fig:regression_r2}). Wniosek
ten jest pozytywny w~sensie metodologicznym: badanie precyzyjnie zlokalizowało
właściwą dźwignię optymalizacji, wskazując, że wysiłek należy kierować na wybór
modelu i~modyfikację jego wag, a~nie na strojenie parametrów generacji.

Po trzecie, kwantyzacja wag okazała się skuteczną i~bezpieczną metodą
przyspieszenia generacji. Zaprojektowana metodyka \emph{warm-cache} pozwoliła
wyizolować \emph{generation\_ms} spod maskującego wpływu prefill (TTFT)
i~zmierzyć rzeczywisty efekt precyzji wag. Wariant Q4\_K\_M zapewnił
$1{,}46\times$ przyspieszenie generacji dla \texttt{gemma3-4b} i~$1{,}67\times$
dla \texttt{gemma4-e2b} względem pełnej precyzji Q8\_0, przy jednoczesnym
zmniejszeniu pliku GGUF o~ok.~40\% i~bez degradacji jakości semantycznej
(tabela~\ref{tab:kwantyzacja_wyniki}). Q4\_K\_M jest zatem optymalnym
kompromisem między rozmiarem, szybkością a~jakością dla modeli rodziny Gemma.

Po czwarte, i~jest to wynik kluczowy dla użyteczności prototypu, udało się
radykalnie zredukować czas do pierwszego tokenu (TTFT). Na Raspberry~Pi~5
pracującym wyłącznie na procesorze (CPU-only) prefill długiego promptu
systemowego dominuje całkowite opóźnienie --- przy zimnym cache TTFT stanowi
ok.~87\% E2E. Utrzymanie modelu oraz pamięci podręcznej kluczy i~wartości
(KV-cache) w~,,ciepłej'' pamięci RAM (\texttt{keep\_alive=30m}) w~połączeniu
z~konstrukcją promptu ze stałym prefiksem (\emph{prefix caching}) skróciło
medianę TTFT z~ok.~27\,387~ms do~ok.~2378~ms, co odpowiada redukcji rzędu
91\% (ostrożnie szacowanej na ok.~84\% w~warunkach produkcyjnych). Ponieważ
urządzenie brzegowe nie dysponuje akceleratorem GPU, to właśnie zachowanie
gorącego stanu w~RAM eliminuje najkosztowniejszy etap i~czyni interaktywnego
asystenta głosowego na CPU-only realnym.

Po piąte, połączenie powyższych dźwigni prowadzi do konkretnej, rekomendowanej
konfiguracji produkcyjnej: \texttt{gemma3:4b} z~kwantyzacją Q4\_K\_M
i~\emph{prefix cachingiem}. Osiąga ona E2E na poziomie ok.~8{,}3~s przy
wskaźniku halucynacji 1{,}17\% --- czternastokrotnie niższym niż druga
najlepsza opcja (16{,}5\%), przy czym różnice są istotne statystycznie
(rozłączne przedziały ufności~95\%, tabela~\ref{tab:wyniki_finalne}).

\subsection{Ocena prototypu względem założeń}

Tezę pracy sformułowano w~dwóch wymiarach (sekcja~\ref{sec:modele} oraz wstęp),
które posłużą tutaj jako punkt odniesienia dla oceny prototypu.

Pierwszy wymiar --- zwiększenie dostępności informacji o~stanie domku
wypoczynkowego dla użytkownika nietechnicznego poprzez lokalny interfejs
głosowy --- został zrealizowany. Zbudowany prototyp o~architekturze
Teacher--Student działa w~kompletnym łańcuchu przetwarzania: od detekcji
aktywności głosowej (VAD) i~transkrypcji mowy (STT), przez dynamiczne budowanie
kontekstu z~danych sensorycznych IoT przez model Nauczyciela, generację
odpowiedzi przez SLM Ucznia na Raspberry~Pi, aż po syntezę mowy (TTS).
Rozdzielenie kosztownej analizy kontekstu (Nauczyciel) od wrażliwej na
opóźnienia obsługi zapytań (Uczeń) sprawia, że czas odpowiedzi asystenta nie
zależy od czasu generacji kontekstu. Poprawność całego łańcucha zweryfikowano
dla scenariuszy odpowiadających kategoriom złotego zbioru testowego, w~tym
przypadków z~brakującymi i~nieaktualnymi danymi.

Drugi wymiar --- wykonalność rozwiązania na platformie sprzętowej o~ograniczonych
zasobach --- oceniono względem trzech założeń przyjętych we wstępie:
\begin{itemize}
  \item \textbf{Krótki czas odpowiedzi:} założenie spełnione. Dzięki kwantyzacji
        Q4\_K\_M oraz utrzymaniu gorącego stanu modelu w~RAM (\emph{prefix
        caching}) osiągnięto E2E na poziomie ok.~8{,}3~s (TTFT ok.~2{,}4~s,
        generacja ok.~5{,}9~s) --- wartość akceptowalną dla asystenta głosowego,
        zwłaszcza w~połączeniu ze strumieniowaniem odpowiedzi i~równoległą
        syntezą mowy.
  \item \textbf{Ograniczenie halucynacji:} założenie spełnione częściowo.
        Rekomendowana konfiguracja \texttt{gemma3:4b} osiąga zaledwie 1{,}17\%
        halucynacji ogółem, jednak w~najtrudniejszych kategoriach
        \texttt{missing\_data} (39{,}4\%) i~\texttt{stale\_data} (29{,}4\%)
        problem pozostaje otwarty i~wymaga interwencji na poziomie wag modelu.
  \item \textbf{Poprawność językowa i~merytoryczna:} założenie spełnione.
        Metryka \emph{polish\_quality} była rejestrowana przez cały cykl badań,
        a~model \texttt{gemma3:4b} zapewnia komunikatywną jakość języka polskiego;
        wysoka \emph{faithfulness} potwierdza zgodność odpowiedzi z~dostarczonym
        kontekstem.
\end{itemize}

Dla rzetelności należy wskazać ograniczenia prototypu i~samych badań.
Porównanie redukcji TTFT oparto na dwóch różnych protokołach (zmienny prompt
przy zimnym cache vs stały prompt przy ciepłym cache), dlatego traktowane jest
jako obserwacja porównawcza z~zastrzeżeniem metodologicznym, a~nie jako
bezpośredni eksperyment cold~vs~warm. Ponadto halucynacje w~kategoriach
\texttt{missing\_data} i~\texttt{stale\_data} nie zostały wyeliminowane przez
żaden z~badanych mechanizmów (strojenie hiperparametrów, kwantyzację, \emph{prefix
caching}), co wyznacza naturalny kierunek dalszych prac.

\subsection{Perspektywy dalszego rozwoju}

Wyniki przeprowadzonych badań wyznaczają kilka naturalnych kierunków dalszego
rozwoju.

Najważniejszym z~nich jest \textbf{fine-tuning modelu docelowego}. Ponieważ
halucynacje w~kategoriach \texttt{missing\_data} i~\texttt{stale\_data} są
własnością wag modelu, a~nie konfiguracji inferencji, jedyną skuteczną metodą
ich redukcji jest modyfikacja wag. Zaplanowany eksperyment (sekcja~\ref{subsec:finetuning_plan})
zakłada dostrojenie metodą LoRA za pomocą biblioteki \texttt{mlx-lm} na~MacBooku~M4,
na zbiorze ok.~400~przykładów w~formacie \texttt{messages}, pokrywających
wszystkie kategorie zachowań docelowych. Pozwoli to zweryfikować hipotezę, czy
model dostrojony --- bez jawnych reguł w~prompcie systemowym --- przewyższa
model bazowy z~regułami na najtrudniejszych scenariuszach.

Drugim kierunkiem jest \textbf{domknięcie zastrzeżenia metodologicznego}
dotyczącego redukcji TTFT: przeprowadzenie bezpośredniego eksperymentu cold~vs~warm
dla identycznego promptu produkcyjnego, który pozwoliłby precyzyjnie skwantyfikować
zysk z~\emph{prefix cachingu} bez porównywania różnych protokołów pomiarowych.

Trzecim jest \textbf{wdrożenie produkcyjne} w~docelowym środowisku Ośrodka
Wypoczynkowego Politechniki Warszawskiej w~Wildze (projekt Wilga~MiLL). Wymaga
ono zastąpienia symulacji odczytów z~pliku CSV strumieniem danych na żywo
(\emph{live feed}) z~fizycznych czujników oraz integracji z~istniejącą
infrastrukturą IoT.

Dalsze kierunki obejmują rozszerzenie puli badanych modeli i~wariantów
kwantyzacji, ewaluację prototypu z~udziałem docelowych użytkowników
nietechnicznych (badania użyteczności) oraz dalszą optymalizację warstw STT/TTS
i~VAD pod kątem opóźnienia i~naturalności interakcji głosowej.